\SourceLeaf{480}

\Needspace{5\baselineskip}\begin{center}ANHANG.\end{center}
\BookRule{0.15\linewidth}

\Needspace{5\baselineskip}\begin{center}\small HANDSCHRIFTLICHE AUFZEICHNUNGEN VON GAUSS.\end{center}


Zu Artikel 40. \textit{Kommt eine dritte Zahl \(C\) hinzu, so sei \(\lambda'\) der größte gemeinsame Teiler der Zahlen \(\lambda,C\), und man bestimme Zahlen \(k,\gamma\) so, dass \(k\lambda+\gamma C=\lambda'\) ist; daraus folgt \(k\alpha A+k\beta B+\gamma C=\lambda'\). Offenbar ist aber \(\lambda'\) ein gemeinsamer Teiler der Zahlen \(A,B.C\), und zwar der größte; denn gäbe es einen größeren \(=\theta\), so wäre \[k\alpha.\frac{A}{\theta}+k\beta.\frac{B}{\theta}+\gamma.\frac{C}{\theta}=\frac{\lambda'}{\theta}\quad\textit{ganze Zahl, Q. E. A.}\]}


\textit{Damit ist das Gesuchte erreicht, wenn man \(k\alpha=a,\ k\beta=b,\ \gamma=c,\ \lambda'=\mu\) setzt.}


Zu Artikel 114. \textit{Der Beweis lässt sich eleganter so führen: \((a^{3n}-a^n)^2=2+(a^{4n}+1)(a^{2n}-2)\), \((a^{3n}+a^n)^2=-2+(a^{4n}+1)(a^{2n}+2)\), und daher \(\surd2\equiv\pm(a^{3n}-a^n),\ \surd-2\equiv\pm(a^{3n}+a^n)\ (\mathrm{mod.}\ 8n+1)\)}


Zu den in Artikel 256 VI angegebenen 16 positiven Determinanten der Form \(8n+5\), für die die Anzahl der eigentlich primitiven Klassen dreimal so groß ist wie die der uneigentlich primitiven, ‘37, .. 573’, sind hinzugefügt: \textit{677, 701, 709, 757, 781, 813, 829, 877, 885, 901, 909, 925, 933, 973, 997; somit sind es unter 125 insgesamt 31.}


Zu den Worten des Artikels 301: Auf diese Weise findet man für die Summe der Anzahlen der Geschlechter für Determinanten \(-\): bis \(-100\) das Ergebnis \(=234,4\), während sie tatsächlich 233 beträgt. — — \textit{Von \(-1\) bis \(-3000\), Tabelle 11166, Formel 11167,9.}


Zu Artikel 336. \textit{Wären alle Zahlen der Form \(2^{2^m}+1\) Primzahlen, so wäre ein hinreichend angenäherter Ausdruck für die Anzahl der hier betrachteten Zahlen \((N)\), die kleiner als die gegebene Zahl \(M\) sind, der folgende: \(\frac12\left(\frac{\log M}{\log2}\right)^2\).}


Zu dem Satz in Artikel 42 über die Teiler eines Polynoms mit ganzzahligen Koeffizienten — — \textit{1797, 22. Juli.}


Zu den Worten des Artikels 130: Nachdem wir streng bewiesen haben, dass jede Primzahl der Form \(4n+1\), sowohl positiv als auch negativ genommen, Nichtrest einer kleineren Primzahl ist, — — \textit{Diesen Beweis entdeckten wir 1796, am 8. April.}


Zu Artikel 131. \textit{Das Fundamentaltheorem wurde im März 1795 durch Induktion entdeckt. Der erste Beweis, der in diesem Abschnitt mitgeteilt wird, wurde im April 1796 gefunden.}

\SourceLeaf{481}


Zu den Worten des Artikels 133: Gehen wir bei der Untersuchung (des Fundamentaltheorems) noch allgemeiner vor. Betrachten wir zwei beliebige ungerade, zueinander teilerfremde Zahlen mit beliebigen Vorzeichen, \(P\) und \(Q\). — — \textit{1796, 29. April.}


Zu den Worten des Artikels 145: Außerdem hätten damals die Sätze über die Reste \(+2\) und \(-2\) vorausgesetzt werden müssen; da unser Beweis jedoch ohne diese Sätze vollständig geführt ist, gewinnen wir daraus eine neue Methode, sie zu beweisen. — — \textit{1797, 4. Februar.}


Zur Überschrift des fünften Abschnitts: Formen und unbestimmte Gleichungen zweiten Grades. — — \textit{Seit dem 22. Juni 1796.}


Zu den Worten des Artikels 234: .. wir gehen zu einem anderen höchst wichtigen, bisher von niemandem berührten Gegenstand über, der Komposition der Formen. — — \textit{Diese Untersuchungen wurden im Herbst 1798 begonnen.}


Zu den Worten des Artikels 262: Aus diesem Prinzip können wir eine neue Methode gewinnen, nicht nur das Fundamentaltheorem, sondern auch die übrigen Sätze des vorigen Abschnitts über die Reste \(-1,\allowbreak +2,\allowbreak -2\) zu beweisen, — — \textit{Die Grundlagen dieser Methode zeigten sich zuerst 1796, am 27. Juli; ausgearbeitet und in ihre gegenwärtige Form gebracht wurden sie aber im Frühjahr 1800.}


Zu den Worten des Artikels 266: .. da aber noch mehrere hierher gehörige, und zwar besonders schöne Wahrheiten übrig bleiben, deren eigentliche Quelle in der Theorie der ternären Formen zweiten Grades zu suchen ist, schalten wir hier einen kurzen Exkurs über diese Theorie ein, — — \textit{14. Februar 1799.}


Zu den Worten des Artikels 272: nämlich indem wir zuerst zeigen, wie sich jede ternäre Form auf eine einfachere Form zurückführen lässt, und dann, dass die Anzahl der einfachsten Formen (zu denen man durch solche Reduktionen gelangt) für jede gegebene Determinante endlich ist. — — \textit{1800, 13. Februar.}


Zu den Worten des Artikels 287 III: Auf ganz ähnliche Weise wird bewiesen, dass in der uneigentlich primitiven Ordnung alle Charaktere, die nach den Vorschriften des Artikels 264 II, III allein als möglich erkannt werden, möglich sind, seien sie \(P\) oder \(Q\). — Diese Sätze usw. — — \textit{Sie wurden erstmals im April 1798 mit Beweisen versehen.}


Zu den Worten des Artikels 302: Die durchschnittliche Anzahl der Klassen aber (die keiner Definition bedürfen wird) wächst sehr regelmäßig, — — \textit{Die erste Idee zu Beginn des Jahres 1799.}


Zu den Worten des Artikels 306 X: Schließlich bemerken wir: Da alle in diesem und im vorigen Artikel betrachteten Eigenschaften vor allem von der Zahl \(n\) abhängen, die etwas Ähnliches ist wie \(p-1\) im dritten Abschnitt, verdient diese Zahl die größte Aufmerksamkeit; daher wäre es höchst wünschenswert, einen allgemeinen Zusammenhang zwischen ihr und der zugehörigen Determinante zu entdecken. — — \textit{Es ist uns wunschgemäß—und so vollständig—gelungen, dass nichts mehr zu wünschen übrig bleibt. 30. November—3. Dezember 1800.}


Zu Artikel 365. \textit{Dass sich der Kreis geometrisch in 17 Teile teilen lässt, entdeckten wir 1796, am 30. März.}

\BookRule{0.3\linewidth}

\clearpage
